\scanpage{38}
\mspage{36}
% archive sheet number: 37
Dans le cas d'un corps de base $k_0$ quelconque, $k$ \deleted{\unsure{$S$ est li}} est limite inductive de ses sous-extensions $k_\alpha$ de t.f. sur le corps premier, et pour $\alpha$ assez grand $\unsure{S_0|S}$ provient d'un $S_\alpha$ sur $k_\alpha$. Soit $x_\alpha$ l'image de $x\in S$ dans $S_\alpha$. Comme $\unsure{\mathscr O_{S,x}|\mathcal O_{S,x}}$ est \uline{plat} sur $\mathscr O_{S_\alpha,x_\alpha}$, on a $\dim\mathscr O_{S_\alpha,x_\alpha}\leq\dim\mathscr O_{S,x}=1$, \uline{et} $\mathscr O_{S_\alpha,x_\alpha}$ est régulier, donc $\mathscr O_{S_\alpha,x_\alpha}$ est un corps ou \deleted{\unsure{un anneau de}} un anneau de val. discrète.
D'ailleurs pour $\alpha$ assez grand, c'est un anneau de valuation discrète, \deleted{et si $\unsure{\mathfrak m_\alpha|m_\alpha}$ est} (sinon $\mathscr O_{S,s}$ serait un corps!), \deleted{de plus $\mathfrak m_\alpha$ est engendré par $t$} et $t\in\mathfrak m_\alpha$. \deleted{Comme \unsure{$\mathfrak m/\mathfrak m^2=\lim\mathfrak m_\alpha/\mathfrak m_\alpha^2$}, on voit que} $t\notin\mathfrak m_\alpha^2$ (sinon on aurait $t\in\mathfrak m^2$), donc on est sous les conditions de a), b) de 8.1.
